\documentclass[12pt]{article}
\usepackage[utf8]{inputenc}
\usepackage[T1]{fontenc}
\usepackage[margin=1in]{geometry}
\usepackage{setspace}
\usepackage{xcolor}
\usepackage[colorlinks=true, urlcolor=blue!60!black, linkcolor=blue!60!black]{hyperref}
\setlength{\parindent}{2em}
\setlength{\parskip}{1.25em}
\renewcommand{\baselinestretch}{1.0}
\begin{document}
\begin{center}
{\LARGE \textit{Carson v. Makin}, 596 U.S. 767 (2022)\par}
\vspace{0.8em}
{\large Case Brief --- Four-Page Edition\par}
\end{center}
\vspace{1em}

\section{Memory Jogger}

Maine's exclusion of ``sectarian'' schools from an otherwise generally available rural tuition-assistance program violates the Free Exercise Clause; there is no safe harbor for discrimination framed around religious ``use'' rather than religious ``status'' (\href{https://www.courtlistener.com/opinion/6479298/carson-v-makin/}{CourtListener}).

\section{Detailed Case Facts}

Maine law guarantees every child a free public education, but because Maine is the most rural State, fewer than half of its 260 school administrative units (SAUs) operate public secondary schools. An SAU that neither runs a high school nor contracts with one must pay tuition, up to a cap, at the public or approved private school the parent selects. Approved private schools may be out of state, single-sex, and selective, need not hire certified teachers, and, if accredited, are exempt from public-school curricular mandates.

Before 1981, parents could choose religious schools. After a 1980 Attorney General opinion asserted that this violated the Establishment Clause, Maine limited the program to ``nonsectarian'' schools, and it kept the restriction even after \textit{Zelman v. Simmons-Harris} (2002) upheld neutral school-choice programs. The Carsons sent their daughter to Bangor Christian Schools at their own expense; the Nelsons could not afford Temple Academy for both children. Both schools were accredited, integrated Christian doctrine across their curricula, and used faith-based admissions and hiring.

\section{Procedural History}

The families sued Commissioner A. Pender Makin. The District of Maine entered judgment for Maine, following First Circuit precedent. After \textit{Espinoza v. Montana Department of Revenue} (2020), the First Circuit still affirmed: the exclusion turned on religious ``use,'' not ``status,'' and the benefit was the ``rough equivalent'' of a secular public education. The Supreme Court reversed and remanded on June 21, 2022.

\section{Judicial Votes}

Reversed, \textbf{6--3}.
\begin{itemize}
\item \textbf{Majority:} Chief Justice Roberts, joined by Justices Thomas, Alito, Gorsuch, Kavanaugh, and Barrett.
\item \textbf{Dissents:} Justice Breyer, joined by Justice Kagan and (except Part I--B) Justice Sotomayor; Justice Sotomayor separately.
\end{itemize}

\section{Holding}

Maine's ``nonsectarian'' requirement violates the Free Exercise Clause. A State need not subsidize private education, but once it offers generally available funding for private schooling, excluding religious schools, whether for religious ``status'' or religious ``use,'' triggers strict scrutiny. A preference for stricter separation than the Establishment Clause requires is not a compelling interest. The Court left open whether neutral antidiscrimination conditions may govern participating schools.

\section{Analysis of Opinions}

\textbf{Majority (Roberts).} The case applied \textit{Trinity Lutheran} and \textit{Espinoza}: the Free Exercise Clause bars benefit exclusions as well as prohibitions. Because parents direct the funds, including religious schools does not offend the Establishment Clause under \textit{Zelman}, so Maine's interest in stricter separation cannot be compelling. Maine was not funding a public-education substitute: participating private schools may reject applicants, exceed the tuition cap, and set their own curricula, and relabeling the benefit as ``secular private education'' would gut \textit{Espinoza}. Maine remained free to expand public options instead. Nor does use-based discrimination fare better: inculcating faith is at the core of a religious school's mission, and auditing how pervasively a school teaches faith would breed entanglement. \textit{Locke v. Davey} was confined to vocational clergy training.

\textbf{Breyer (dissenting).} The majority ignored the Establishment Clause and abandoned the ``play in the joints'' between the Clauses, converting what \textit{Zelman} permits into what the Constitution requires. These schools wove Christian doctrine through every subject, and Maine's exclusion of religious indoctrination from a civic public-education program fell within \textit{Locke}. In Part I--B, joined only by Justice Kagan, he stressed the Religion Clauses' shared purpose of preventing religious strife in a nation of more than 100 religious groups, invoking Jefferson and Madison against compelled support for religious teaching.

\textbf{Sotomayor (dissenting).} Her \textit{Trinity Lutheran} warning that status-based holdings would erase use-based limits had come true; Maine taxpayers must now fund schools that exclude students and teachers on the basis of religion, sexual orientation, and gender identity. Maine's program, she added, was driven by rural geography, not an elective preference for private subsidies, and the Court had made separation of church and state itself a constitutional violation.

\section{Examples: Future Applications}

\begin{enumerate}
\item \textbf{Same-Side: Education Savings Accounts.} A State's \$7,000 accounts bar schools that integrate religion into academic courses. Unconstitutional: a use-based exclusion from a parent-directed benefit fails strict scrutiny.
\item \textbf{Same-Side: Use-Based Blaine Amendment.} Officials invoke a state ban on funds ``applied to any religious worship, exercise, or instruction'' to exclude religious schools from a tax-credit scholarship. Unconstitutional after \textit{Carson}.
\item \textbf{Opposite-Side: Purely Public Delivery.} A rural State ends tuition payments and builds regional public schools, busing, boarding dormitories, and online classes. Constitutional: a State need not fund private schooling at all.
\item \textbf{Opposite-Side: Divinity Degrees.} A scholarship excludes students training for the clergy. Constitutional under \textit{Locke v. Davey}, which \textit{Carson} preserved.
\item \textbf{Fence-Sitter: Antidiscrimination Conditions.} A State drops its ``nonsectarian'' clause but requires all participating schools to follow nondiscrimination laws in admissions and hiring. Footnote 1* left this open, pitting \textit{Employment Division v. Smith} against \textit{Fulton} and the ministerial exception. A related question, whether a State must charter a religious charter school, later deadlocked the Court 4--4 in \textit{Oklahoma Statewide Charter School Board v. Drummond} (2025).
\end{enumerate}

\section{Critique}

\textbf{Progressive critique.} Ira C. Lupu and Robert W. Tuttle argue that \textit{Carson} swept aside decades of precedent forbidding taxpayer financing of devotional instruction, discarding the status/use distinction when \textit{Zelman} offered a narrower ground and threatening ``Charitable Choice'' rules; they note the asymmetry of exempting religious schools from employment law while compelling taxpayers to fund them. Critics further warned that pressuring States either to finance schools that discriminate against LGBTQ students or to eliminate rural tuition programs harms both public education and religious minorities.

\textbf{Originalist and school-choice defense.} Supporters view \textit{Carson} as the coherent completion of \textit{Zelman}, \textit{Trinity Lutheran}, and \textit{Espinoza}. Michael Bindas and others praise \textit{Carson} for protecting religious \textit{exercise}, not mere identity, and for ending intrusive theological line-drawing; Blaine Amendments reflect nineteenth-century anti-Catholic nativism. Some federalists still regret shrinking States' ``play in the joints.''

\section{Key Quotations}

\begin{quote}
``In short, the prohibition on status-based discrimination under the Free Exercise Clause is not a permission to engage in use-based discrimination.''
\end{quote}
\noindent --- Chief Justice Roberts, for the Court, 596 U.S. 767.

\begin{quote}
``The State pays tuition for certain students at private schools---so long as the schools are not religious. That is discrimination against religion.''
\end{quote}
\noindent --- Chief Justice Roberts, for the Court, 596 U.S. 767.

\begin{quote}
``Regardless of how the benefit and restriction are described, the program operates to identify and exclude otherwise eligible schools on the basis of their religious exercise.''
\end{quote}
\noindent --- Chief Justice Roberts, for the Court, 596 U.S. 767.

\end{document}
